\documentclass{article}
\usepackage[utf8]{inputenc}
\usepackage[english, spanish]{babel}
%\usepackage[dvips]{graphics}
\usepackage{amsmath}
\usepackage{amssymb}
\usepackage{fullpage}
\usepackage{epsfig}
\usepackage{multicol}
\usepackage{wasysym}
\usepackage{tikz}
\usepackage{graphicx}
\usepackage{mathtools}
\usepackage{enumitem}

\usetikzlibrary{arrows,decorations.pathmorphing,backgrounds,positioning,fit,calc,automata}
\usetikzlibrary{decorations.pathreplacing}
\usetikzlibrary{trees}
\usetikzlibrary{shapes}
\usetikzlibrary{chains}
\usetikzlibrary{patterns}
\usetikzlibrary{matrix}
%\usepackage{stackengine}


%\usepackage[usenames,dvipsnames]{xcolor}
%\usepackage{hyperref}
%\definecolor{linkcolour}{rgb}{0,0.2,0.6}
%\hypersetup{colorlinks,breaklinks,urlcolor=linkcolour,linkcolor=linkcolour}

\parindent 0pt
%\parskip 0pt

\newcommand{\cL}{\mathcal{L}}

\begin{document}

\includegraphics[width=2cm]{uc.png}
\vspace*{-1.9cm}

\hspace*{2.1cm}
 \begin{tabular}{l}
  \sc Pontificia Universidad Católica de Chile \\
  \sc Escuela de Ingeniería \\
  \sc Departamento de Ciencia de la Computación
  %\vspace{15\baselineskip}\mbox{}
  %\vspace{-3mm}\mbox{}
 \end{tabular}
 \bigskip

\vspace*{5mm}
\begin{center}
{IIC2213 --- Lógica para ciencia de la computación --- 1' 2023} \\
\vspace{3mm}
{\Large\bf TAREA 1} \\
\vspace{2mm}
\end{center}

\begin{tabular}{ll}
Publicación: & Martes 14 de marzo. \\
Entrega: & \textbf{Lunes 27 de marzo hasta las 23:59 horas}. \\
\end{tabular}

\subsection*{Indicaciones}

\begin{itemize}
\item Cada pregunta tiene 6 puntos ($+1$ base) y la nota de la tarea es el promedio de  las preguntas.
\item Todos los incisos de una pregunta pregunta tienen el mismo puntaje.
\item Cada solución debe estar escrita en \LaTeX. No se aceptarán tareas escritas de otra forma.
\item La tarea es individual, pudiendo discutirla con sus pares. Toda referencia externa debe citarse.
\end{itemize}

\subsection*{Objetivos}
Esta tarea busca evaluar si el estudiante es capaz de
\begin{itemize}
    \item Definir inductivamente propiedades/objetos en $\cL(P)$
    \item Demostrar usando diversas técnicas, entre ellas inducción en $\cL(P)$
    \item Aplicar conceptos de sintaxis y semántica proposicional, y consecuencia lógica
\end{itemize}

\subsection*{Pregunta 1: Sintaxis e inducción en $\cL(P)$}
Decimos que dos grafos $G_1 = (V_1, E_1)$ y $G_2 =
(V_2, E_2)$\ son {\em isomorfos}, denotado por $G_1\cong G_2$, si existe una biyección $f : V_1
\rightarrow V_2$ tal que para todo $a,b\in V_1$ se tiene que
$$\{a,b\} \in E_1\quad\text{si y solo si}\quad\{f(a), f(b)\} \in E_2$$ 

Dado un par de grafos no dirigidos $G_1=(V_1,E_1)$ y $G_2=(V_2,E_2)$, construya una fórmula proposicional $\varphi_\text{iso}$ tal que 
$$\varphi_\text{iso}\text{ es satisfacible\quad si y solo si}\quad G_1\cong G_2$$
Demuestre la correctitud de su construcción.

\subsubsection*{Solución P1.}
Aquí va mi solución

\newpage
\subsection*{Pregunta 2: Semántica en $\cL(P)$ y equivalencia lógica}
Explique un procedimiento que utilice resolución para determinar si una fórmula $\varphi$ cualquiera es tautología. Use su método para determinar si
$\varphi=(p\rightarrow q)\rightarrow(\neg p\rightarrow\neg q)$
es tautología.

\subsubsection*{Solución P2.}
Aquí va mi solución

\newpage
\subsection*{Pregunta 3: Consecuencia lógica en $\cL(P)$}
\begin{enumerate}
    \item[(a)] Sea $\mathcal{L}$ un vocabulario arbitrario. Demuestre que si una propiedad $\mathcal{P}$ es elementalmente definible, entonces la propiedad
    \begin{center}
        $\overline{\mathcal{P}}=\textsc{S}[\mathcal{L}]\setminus\mathcal{P}$
    \end{center}
    también es elementalmente definible.
    \item[(b)] Considere el vocabulario $\cL=\{a,b,E\}$ con símbolos de constantes $a$ y $b$, y símbolo de relación $E$ binaria. Demuestre usando el teorema de compacidad que la siguiente propiedad no es elementalmente definible
    \begin{center}
        $\mathcal{P}=\{\mathfrak{A}\in\textsc{S}[\mathcal{L}]\mid \mathfrak{A}\text{ tiene un camino de largo finito entre }a^\mathfrak{A}\text{ y }b^\mathfrak{A}\}$
    \end{center}
\end{enumerate}

\subsubsection*{Solución P3.}
Aquí va mi solución


\end{document}


